% ID: SEL-D0019-OB-E02
% ESTADO: cotejada. Segunda pasada visual completa y borrador compilado revisados.
% Fuente: Selectividad/matematicas_aplicadas/reserva a/Examen-Matemáticas Aplicadas a las Ciencias Sociales.pdf
% SHA-256: 6fed0384ff6fcbaaa7098ab3f35ea5f881d1280502f050ada3a7f13af528392e
% Página: 2; opción B; ejercicio original 2.

Se considera la función $f(x)=\dfrac{ax}{bx+1}$, con $a$ y $b$ números reales.
\begin{enumerate}[label=\alph*)]
\item \textbf{(1.5 puntos)} Calcule los valores de $a$ y $b$, sabiendo que $f(-1)=1$ y que en el punto de abscisa $x=0$ la recta tangente a la gráfica de $f$ es paralela a la recta $y=2x+1$.
\item \textbf{(1 punto)} Para $a=b=1$, halle la ecuación de sus asíntotas.
\end{enumerate}
